\section{Language, calculus, and Kripke structures}
\label{sec:intuitionistic-logic-01}

The propositional language contains
\(\land,\lor,\to,\bottom\), and
\[
  \neg A\defeq A\to\bottom.
\]
Because \(\LJ\)-antecedents are finite multisets whereas premise collections
may be arbitrary sets, derivability is defined by finite support:
\[
  \Gamma\derives_{\IPL}A
  \quad:\Longleftrightarrow\quad
  \begin{gathered}
    \text{there is a finite multiset context \(\Gamma_0\), every formula}\\
    \text{occurring in which belongs to \(\Gamma\), such that}
    \quad \LJ\derives\Gamma_0\Rightarrow A.
  \end{gathered}
\]
This defines the left side for every formula set \(\Gamma\) while using only
contexts admitted by \(\LJ\).  The unlifted sequent notation remains available
for a particular finite multiset \(\Gamma_0\).

\begin{definition}[Intuitionistic Kripke frame and structure]
An intuitionistic Kripke frame is a partially ordered set \((W,\leq)\).  A
Kripke structure is \(\mathcal M=(W,\leq,V)\), where
\(V:\Prop\to\Pow(W)\) is persistent:
\[
  w\in V(p)\text{ and }w\leq v
  \quad\Longrightarrow\quad
  v\in V(p).
\]
The order represents growth of information.  This interpretation motivates
the semantics but is not part of its formal definition.
\end{definition}

The forcing relation \(\mathcal M,w\forces A\) is defined recursively:
\[
\begin{aligned}
  \mathcal M,w\forces p
    &\quad\Longleftrightarrow\quad w\in V(p),\\
  \mathcal M,w\not\forces\bottom,
    &\\
  \mathcal M,w\forces A\land B
    &\quad\Longleftrightarrow\quad
      \mathcal M,w\forces A\text{ and }\mathcal M,w\forces B,\\
  \mathcal M,w\forces A\lor B
    &\quad\Longleftrightarrow\quad
      \mathcal M,w\forces A\text{ or }\mathcal M,w\forces B,\\
  \mathcal M,w\forces A\to B
    &\quad\Longleftrightarrow\quad
      \forall v\geq w\,
      (\mathcal M,v\forces A\to\mathcal M,v\forces B).
\end{aligned}
\]
The arrow inside the final parenthesis is metalinguistic implication.

